\begin{thebibliography}{99}

\bibitem{biedenharnLouck1981}
L. C. Biedenharn and J. D. Louck,
\emph{Angular Momentum in Quantum Physics: Theory and Application},
Encyclopedia of Mathematics and its Applications, Vol.~8,
Addison--Wesley, 1981.

\bibitem{kirillovReshetikhin1989}
A. N. Kirillov and N. Yu. Reshetikhin,
``Representations of the algebra $U_q(\mathfrak{sl}_2)$, q-orthogonal polynomials and invariants of links,''
in \emph{Infinite-Dimensional Lie Algebras and Groups},
Advanced Series in Mathematical Physics, Vol.~7,
World Scientific, 1989, pp.~285--339.

\bibitem{kauffmanLins1994}
L. H. Kauffman and S. L. Lins,
\emph{Temperley--Lieb Recoupling Theory and Invariants of 3-Manifolds},
Annals of Mathematics Studies, Vol.~134,
Princeton University Press, 1994.

\bibitem{ponzanoRegge1968}
G. Ponzano and T. Regge,
``Semiclassical limit of Racah coefficients,''
in \emph{Spectroscopic and Group Theoretical Methods in Physics},
North-Holland, Amsterdam, 1968, pp.~1--58.

\bibitem{turaevViro1992}
V. G. Turaev and O. Ya. Viro,
``State sum invariants of 3-manifolds and quantum $6j$-symbols,''
\emph{Topology} \textbf{31} (1992), no.~4, 865--902.
\href{https://doi.org/10.1016/0040-9383(92)90015-A}{doi:10.1016/0040-9383(92)90015-A}.

\bibitem{roberts1999}
J. Roberts,
``Classical $6j$-symbols and the tetrahedron,''
\emph{Geometry \& Topology} \textbf{3} (1999), 21--66.
\href{https://doi.org/10.2140/gt.1999.3.21}{doi:10.2140/gt.1999.3.21};
\href{https://arxiv.org/abs/math-ph/9812013}{arXiv:math-ph/9812013}.

\bibitem{taylorWoodward2005}
Y. U. Taylor and C. T. Woodward,
``$6j$ symbols for $U_q(\mathfrak{sl}_2)$ and non-Euclidean tetrahedra,''
\emph{Selecta Mathematica} \textbf{11} (2005), 539--571.
\href{https://arxiv.org/abs/math/0305113}{arXiv:math/0305113}.

\bibitem{kato1995}
T. Kato,
\emph{Perturbation Theory for Linear Operators}, 2nd ed.,
Springer, 1995.

\bibitem{ahlfors1979}
L. V. Ahlfors,
\emph{Complex Analysis}, 3rd ed.,
McGraw--Hill, 1979.

\bibitem{johansson2017arb}
F. Johansson,
``Arb: efficient arbitrary-precision midpoint-radius interval arithmetic,''
\emph{IEEE Transactions on Computers} \textbf{66} (2017), no.~8, 1281--1292.
\href{https://doi.org/10.1109/TC.2017.2690633}{doi:10.1109/TC.2017.2690633}.

\end{thebibliography}
